\begin{abstract}
\noindent
Dieser Bericht dokumentiert das Forschungsprojekt B. Im Zentrum steht die Frage, ob sich mithilfe Hessian-basierter Krümmungsmetriken vorhersagen lässt, welche Schichten eines Convolutional Neural Network unter der Power-of-Two-Quantisierung (PoT) am stärksten geschaedigt werden, wie Post-Training-Quantisierung und Quantization-Aware-Training diese Krümmung dabei jeweils unterschiedlich prägen, und in welchem Ausmaß der Schaden überhaupt auf Gewichts- statt Aktivierungsquantisierung zurückgeht. Die PoT ist ein hardware-attraktives Quantisierungsschema, da sie Multiplikationen durch Bitshifts ersetzt. Gleichzeitig erzeugt sie jedoch ungleichmäßige Genauigkeitsverluste über die Schichten hinweg. Als etablierter Ansatz schlägt die HAWQ-V2-Formulierung hierfür das Produkt aus Hessian-Spur und Störungsgröße, $\mathrm{Tr}(H)\cdot\|\Delta W\|^2$, als Sensitivitaets-Score vor. Um den Schaden jeder Schicht isoliert zu bestimmen, wird jeweils nur diese eine Schicht quantisiert und der daraus resultierende Anstieg des Validierungs-Loss gemessen. Auf dieser Grundlage werden drei Scores gegen den gemessenen Schaden getestet, die rohe Krümmung, die Stoerungsgröße und das HAWQ-V2-Produkt, und zwar über drei Architekturen (CNN, ResNet-18, ResNet-50) sowie zwei Datensaetze (CIFAR10, ImageNet100). Die Ergebnisse zeigen, dass kein Score universell zuverlässig ist, und dass auch die rohe Krümmung nur ein eng umrissenes, schwaches Signal liefert: Sie identifiziert die am stärksten geschädigte Schicht (Rang~1) in 3 von 6 Architektur$\times$Datensatz-Kombinationen, doch schließen 5 von 6 Bootstrap-Konfidenzintervallen der Rangkorrelation Null ein; im einzigen signifikanten Fall (ResNet-50/CIFAR10) platziert dieselbe Krümmung die wahre Top-Schicht auf Rang 40 von 54. Zwei der drei Rang-1-Erfolge erweisen sich zudem als direkt an \texttt{conv1} gebunden, das selbst überwiegend ein Fusions-/Fan-in-Artefakt ist, und verschwinden, sobald \texttt{conv1} aus der Schichtmenge entfernt wird. Zwei günstige Baselines ganz ohne Hessian-Berechnung, Gradientennorm und empirische Fisher-Diagonale, erreichen bei einem Bruchteil des Rechenaufwands dieselbe oder eine leicht bessere Rang- und top-$k^*$-Trefferquote wie die rohe Krümmung. Das HAWQ-V2-Produkt wird zusätzlich durch eine Anti-Korrelation zwischen Störungsgröße und tatsaechlichem Schaden in die Irre geführt. Um der Kruemmung auf den Grund zu gehen wurden einige Analyse-Methoden angewandt. Eine Random-Init-Kontrolle, eine gescheiterte KFAC-Zerlegung, eine Dekomposition in Gewichts- und Aktivierungsanteile sowie eine Deployment-Evaluation auf realer Hardware. Dabei zeigt sich, dass die PoT-Quantisierung auf einer NVIDIA A100 keinen Laufzeitvorteil bringt, da aktuelle INT8-Kernel für Faltungsschichten fehlen, während ein CPU-Pfad mit echten INT8-Kernen (fbgemm, CIFAR10) einen realen Durchsatzgewinn misst, allerdings für generische INT8-Arithmetik und nicht für den PoT-spezifischen Bitshift selbst, der auf keiner der beiden Plattformen nativ verfügbar und somit ungemessen bleibt. Ein zentrales methodisches Ergebnis ist zudem ein Rekonziliations-Ledger, das dokumentiert, wie stark Trace-Werte von Konfigurationsentscheidungen abhängen, die in der Literatur meist unausgesprochen bleiben. Die Ergebnisse der Kernfragestellung münden schließlich in ein Workshop-Paper, das bei AXIOM @ NeurIPS 2026 eingereicht wurde.
\end{abstract}